\chapter{RELATED WORK AND BACKGROUND}

\section{Algorithmic Foundations of Shortest-Path Search}

The foundation of this work rests on Dijkstra’s algorithm~\citep{dijkstra1959}, a cornerstone of graph theory for finding the shortest path in non-negative weighted graphs. By maintaining a priority queue to explore nodes in order of increasing accumulated cost, Dijkstra’s algorithm guarantees the identification of the globally optimal path. When implemented with a binary min-heap, the algorithm achieves a time complexity of $O(|E|\log|V|)$~\citep{cormen2009}, which is utilized as our primary baseline. 

In the context of customer journeys, we apply a log-probability transformation $w(u,v) = -\log\,p(u,v)$. This exploits the monotone relationship between probability maximization and additive cost minimization—a standard technique in probabilistic graph optimization~\citep{manning1999nlp}. Adopting the foundations established by \citet{cormen2009}, our implementation utilizes adjacency lists for $O(|V|+|E|)$ space efficiency and hash maps for $O(1)$ distance lookups, ensuring scalability for the large-scale e-commerce graphs analyzed in Chapter 6.

\section{Heuristic vs. Threshold-Based Approximation}

While Dijkstra’s algorithm ensures optimality, its exhaustive nature may lead to unnecessary expansion of low-probability states. Informed search algorithms, notably A*~\citep{hart1968astar}, utilize an admissible heuristic $h(n)$ to guide the search. However, as noted by \citet{pearl1984heuristics}, designing a tight, admissible heuristic for log-probability weights in arbitrary navigation graphs is non-trivial. 

Alternative approximations, such as Beam Search, limit the search frontier to the top-$k$ candidates~\citep{cormen2009}. While efficient, Beam Search lacks a formal convergence guarantee and may discard the optimal path entirely. The proposed \textbf{Probability-Pruned Dijkstra} variant occupies a middle ground. It is a deterministic, parameter-controlled framework that uses a threshold $\tau$ to prune paths that fall below a minimum probability. Unlike rank-based pruning (Beam Search), this method provides a formal \emph{All-or-Nothing Property}: it either returns the globally optimal path or reports no path at all, ensuring that accuracy is never sacrificed for speed when a result is returned.

\section{Probabilistic Navigation and Markovian Assumptions}

Modeling user interactions as a sequence of state transitions is formally equivalent to a first-order Markov chain~\citep{norris1998markov}. This paradigm has been extensively validated in web navigation and clickstream analysis~\citep{bucklin2003clickstream, liu2011webmining}. The transition-based models discussed by \citet{resnick1997recommender} serve as the basis for our path-finding approach.

The maximum-probability path problem is related to the Viterbi algorithm~\citep{viterbi1967error, rabiner1989hmm}, which finds the most likely state sequence in Hidden Markov Models. However, while Viterbi is optimized for fixed-length sequences, our framework generalizes the principle to arbitrary directed graphs with cycles and variable path lengths. We assume a First-Order Markov property where the next state depends only on the current state; this assumption is empirically grounded in our analysis of the RetailRocket and YOOCHOOSE datasets, which exhibit stable transition matrices across high-volume interaction sessions.

\section{Alternative Computational Paradigms}

Customer journey analysis often intersects with \emph{Sequential Pattern Mining} and \emph{Scalable Graph Processing}. Algorithms like PrefixSpan~\citep{pei2001prefixspan} and general surveys of frequent pattern mining~\citep{han2007sequential} focus on discovering frequent sequences rather than optimizing specific source-to-target queries. While these mining techniques utilize pruning (e.g., minimum support thresholds), they serve a descriptive rather than an optimization-based purpose.

Furthermore, scalable systems like Pregel~\citep{malewicz2010pregel} and subsequent distributed graph frameworks~\citep{sahu2017scaling} address the engineering challenges of massive graphs. Our work focuses instead on the \emph{algorithmic characterization} of pruning. By evaluating search-space reduction within a controlled single-machine environment, we establish the fundamental trade-offs between pruning intensity and path reachability that would inform implementations in distributed environments.

\section{Structural Characteristics: Cycles and Sparsity}

A critical structural property of user navigation is the presence of cycles (e.g., a user returning to a "Browse" state from a "Product" page). Following \citet{bucklin2003clickstream}, we model the journey as a general directed cyclic graph. Because our log-transformed weights are strictly non-negative ($w \geq 0$), Dijkstra’s algorithm remains correct and naturally avoids infinite loops, as cycle traversal strictly increases path cost~\citep{cormen2009}.

Additionally, our framework assumes graph sparsity where $|E| \approx O(|V|)$. This assumption is not merely theoretical; it is empirically supported by the datasets used in this study. For instance, the item-level transition graphs derived from RetailRocket exhibit a low average out-degree, reflecting the reality that users typically navigate between a limited set of related products rather than a fully connected network.

\section{Background on Evaluation Data}

To ensure the robustness of the pruning framework, we utilize a dual-pronged evaluation strategy involving both synthetic and real-world data:

\begin{itemize}
    \item \textbf{Synthetic Graphs:} Grounded in the Erd\H{o}s--R\'{e}nyi random graph model~\citep{erdosrenyi1960}, we implement two generators: a sparse random generator and a layered stage-based generator. These allow for controlled stress-testing of $\tau$ across varying node densities and stage-gate structures.
    \item \textbf{Real-World Datasets:} We utilize the \textbf{RetailRocket}~\citep{retailrocket2015} and \textbf{RecSys 2015 YOOCHOOSE}~\citep{recsys2015} datasets. These provide millions of clickstream events to validate the algorithm on actual user behavioral patterns, moving the evaluation beyond purely theoretical or synthetic bounds.
\end{itemize}

\section{Research Gap and Project Scope}

Despite the maturity of shortest-path theory and clickstream modeling, a gap exists in the systematic, empirical evaluation of \emph{threshold-based pruning} specifically for cyclic customer journey graphs. Existing literature tends to treat sequential mining~\citep{pei2001prefixspan} and shortest-path search~\citep{dijkstra1959} as separate domains. 

This project addresses this gap by presenting a **controlled pruning framework** that integrates probabilistic transition modeling with formal optimality characterizations. Rather than proposing a fundamentally new shortest-path theory, this work focuses on the empirical characterization of the time–optimality trade-off, providing a principled method for reducing the search space in high-scale conversion optimization problems while maintaining a mathematical guarantee of path optimality.